% ===== Stack (7 problems) =====
\section{Stack \texorpdfstring{{\normalsize\color{muted}$\cdot$ 7 problems}}{- 7 problems}}
\markboth{Stack}{}

\begin{problem}{20}{Valid Parentheses}{Easy}
\summary{Are brackets ()[]\{\} correctly matched and nested?}
\keyidea{Push openers; on a closer, the top must be its matching opener.}
\complexity{O(n) time, O(n) space}
\begin{code}
class Solution:
    def isValid(self, s: str) -> bool:
        pairs = {")": "(", "]": "[", "}": "{"}
        stack = []
        for c in s:
            if c in pairs:
                if not stack or stack[-1] != pairs[c]:
                    return False
                stack.pop()
            else:
                stack.append(c)
        return not stack
\end{code}
\end{problem}

\begin{problem}{155}{Min Stack}{Medium}
\summary{Stack supporting push, pop, top and getMin, all O(1).}
\keyidea{Second stack storing the min at each depth.}
\complexity{O(1) per op, O(n) space}
\begin{code}
class MinStack:
    def __init__(self):
        self.stack, self.mins = [], []

    def push(self, val: int) -> None:
        self.stack.append(val)
        self.mins.append(min(val, self.mins[-1] if self.mins else val))

    def pop(self) -> None:
        self.stack.pop()
        self.mins.pop()

    def top(self) -> int:
        return self.stack[-1]

    def getMin(self) -> int:
        return self.mins[-1]
\end{code}
\end{problem}

\begin{problem}{150}{Evaluate Reverse Polish Notation}{Medium}
\summary{Evaluate an expression in postfix notation (+ - * /, division truncates toward zero).}
\keyidea{Push numbers; on an operator pop b then a, push a op b. Use int(a / b) for truncation.}
\complexity{O(n) time, O(n) space}
\begin{code}
class Solution:
    def evalRPN(self, tokens: List[str]) -> int:
        st = []
        for t in tokens:
            if t in "+-*/":
                b, a = st.pop(), st.pop()
                if t == "+": st.append(a + b)
                elif t == "-": st.append(a - b)
                elif t == "*": st.append(a * b)
                else: st.append(int(a / b))
            else:
                st.append(int(t))
        return st[0]
\end{code}
\end{problem}

\begin{problem}{22}{Generate Parentheses}{Medium}
\summary{List every well-formed string of n pairs of parentheses.}
\keyidea{Backtrack: add '(' if open \textless{} n; add ')' if close \textless{} open.}
\complexity{O(4\textasciicircum{}n / sqrt(n)) time (Catalan), O(n) recursion}
\begin{code}
class Solution:
    def generateParenthesis(self, n: int) -> List[str]:
        res, cur = [], []
        def bt(op, cl):
            if op == cl == n:
                res.append("".join(cur))
                return
            if op < n:
                cur.append("("); bt(op + 1, cl); cur.pop()
            if cl < op:
                cur.append(")"); bt(op, cl + 1); cur.pop()
        bt(0, 0)
        return res
\end{code}
\end{problem}

\begin{problem}{739}{Daily Temperatures}{Medium}
\summary{For each day, how many days until a warmer temperature (0 if never)?}
\keyidea{Monotonic decreasing stack of indices; a warmer day resolves everything colder on the stack.}
\complexity{O(n) time, O(n) space}
\begin{code}
class Solution:
    def dailyTemperatures(self, temperatures: List[int]) -> List[int]:
        res = [0] * len(temperatures)
        st = []  # indices
        for i, t in enumerate(temperatures):
            while st and temperatures[st[-1]] < t:
                j = st.pop()
                res[j] = i - j
            st.append(i)
        return res
\end{code}
\end{problem}

\begin{problem}{853}{Car Fleet}{Medium}
\summary{Cars drive to target; a faster car that catches a slower one joins it. Count arriving fleets.}
\keyidea{Sort by position descending; compute arrival time. A car forms a new fleet only if it takes longer than the fleet ahead.}
\complexity{O(n log n) time, O(n) space}
\begin{code}
class Solution:
    def carFleet(self, target: int, position: List[int], speed: List[int]) -> int:
        cars = sorted(zip(position, speed), reverse=True)
        st = []
        for p, s in cars:
            t = (target - p) / s
            if not st or t > st[-1]:
                st.append(t)
        return len(st)
\end{code}
\end{problem}

\begin{problem}{84}{Largest Rectangle in Histogram}{Hard}
\summary{Area of the largest rectangle that fits under the histogram bars.}
\keyidea{Increasing stack of (start, height). When popping a taller bar, its rectangle ends here; the new bar inherits the popped start.}
\complexity{O(n) time, O(n) space}
\begin{code}
class Solution:
    def largestRectangleArea(self, heights: List[int]) -> int:
        best, st = 0, []  # (start index, height)
        for i, h in enumerate(heights):
            start = i
            while st and st[-1][1] > h:
                idx, height = st.pop()
                best = max(best, height * (i - idx))
                start = idx
            st.append((start, h))
        for idx, height in st:
            best = max(best, height * (len(heights) - idx))
        return best
\end{code}
\end{problem}
